\section{Ferramentas de indexação}
Escreva $\N_0=\{0,1,2,\ldots\}$ e $T_m=m(m+1)/2$.

\subsection{Pareamento de Cantor}
Para $u,v\in\N_0$, defina
\eqn{01}{\pi(u,v)=T_{u+v}+v=\frac{(u+v)(u+v+1)}2+v.}
\begin{lema}[pareamento e cota]
$\pi:\N_0^2\to\N_0$ é bijetiva. Para $n\in\N_0$, sua inversa é
\eqn{02}{w=\left\lfloor\frac{\sqrt{8n+1}-1}{2}\right\rfloor,\qquad v=n-T_w,\qquad u=w-v.}
Além disso, $\pi(u,v)\le(u+v+1)^2$.
\end{lema}
\begin{proof}
Se $m=u+v$, então $T_m\le\pi(u,v)=T_m+v<T_{m+1}$, pois $0\le v\le m$. Assim, cada imagem pertence a um único bloco diagonal $[T_m,T_{m+1})$. Dado $n$, a expressão para $w$ em \eqref{eq:02} é precisamente o único inteiro com $T_w\le n<T_{w+1}$. Logo $0\le v=n-T_w\le w$, $u=w-v\ge0$ e $\pi(u,v)=n$, provando sobrejetividade. O bloco determina $w=u+v$ e, dentro dele, $v=n-T_w$; portanto o par é único. Por fim, com $x=u+v+1$, tem-se $T_{u+v}\le x^2/2$ e $v\le x^2/2$, donde $\pi(u,v)\le x^2$.
\end{proof}

\subsection{Inteiros e racionais por altura}
Defina $Z:\N_0\to\Z$ por
\eqn{03}{Z(t)=\begin{cases}t/2,&t\text{ par},\\-(t+1)/2,&t\text{ ímpar}.\end{cases}}
\begin{lema}
$Z$ é bijetiva, com $Z^{-1}(m)=2m$ para $m\ge0$ e $Z^{-1}(m)=-2m-1$ para $m<0$. Em particular, $Z^{-1}(m)\le2|m|+1$.
\end{lema}
\begin{proof}
Os índices pares percorrem $\N_0$ exatamente uma vez; os ímpares percorrem $\Z\setminus\N_0$ exatamente uma vez. As duas imagens são disjuntas e as fórmulas inversas resultam por resolução direta de \eqref{eq:03}.
\end{proof}

Todo $r\in\Q$ tem uma representação reduzida única $p/q$, com $p\in\Z$, $q\in\N$ e $\gcd(|p|,q)=1$. Adote a altura multiplicativa elementar
\eqn{04}{\heightq(p/q)=\max\{|p|,q\}.}
Ordene os pares reduzidos primeiro por altura, depois por $q$ e finalmente por $p$. Essa leitura define bijeções
\eqn{05}{Q:\N_0\to\Q,\qquad Q^*:\N_0\to\Q^*,}
na segunda das quais se omite $0/1$.

\begin{lema}[contagem por altura]
Seja $\Q_H=\{r\in\Q:\heightq(r)\le H\}$. Então
\eqn{06}{|\Q_H|\le3H^2\qquad\text{para todo }H\ge1,}
e existem $H_0$ e $c_1>0$ universais tais que
\eqn{07}{|\Q_H|\ge c_1H^2\qquad\text{para todo }H\ge H_0.}
Pode-se tomar $c_1=1/4$.
\end{lema}
\begin{proof}
A forma reduzida com denominador positivo é única: de $p/q=p'/q'$ e das duas coprimalidades segue $q\mid q'$ e $q'\mid q$, logo $q=q'$ e $p=p'$. Cada $r\in\Q_H$ corresponde então a um único par $(p,q)$ com $|p|\le H$ e $1\le q\le H$; há no máximo $(2H+1)H\le3H^2$ tais pares, o que prova \eqref{eq:06}.

Para \eqref{eq:07}, observe que $\Q_H$ contém todas as frações $p/q$ com $1\le q\le H$, $1\le p\le q$ e $\gcd(p,q)=1$, cuja quantidade é $\Phi(H)=\sum_{q\le H}\varphi(q)$. Pelo teorema clássico de Mertens sobre a soma de Euler \citep[Teorema 330]{hardy2008},
\[
\Phi(H)=\frac3{\pi^2}H^2+O(H\log H),
\]
e $3/\pi^2>0{,}30$. Logo $\Phi(H)\ge H^2/4$ para todo $H$ suficientemente grande.
\end{proof}

\begin{lema}[posto por altura]
As aplicações em \eqref{eq:05} são bijetivas. Se $\heightq(p/q)\le H$, então o posto de $p/q$ em qualquer uma delas é estritamente menor que $3H^2$.
\end{lema}
\begin{proof}
A ordenação esgota cada casca finita de altura antes de passar à seguinte; portanto lista todo racional uma única vez, o que dá a bijetividade. Todo elemento que precede $p/q$ pertence a $\Q_H$, e o lema~2.3 limita essa cardinalidade por $3H^2$. A remoção de $0/1$ em $Q^*$ só diminui a contagem.
\end{proof}

\subsection{Bijetividade e crescimento do índice}
A ordem de Calkin--Wilf fornece uma comparação entre bijetividade e crescimento do endereço. Seja $(S_n)_{n\ge0}$ a sequência diatômica de Stern,
\eqn{08}{S_0=0,\qquad S_1=1,\qquad S_{2n}=S_n,\qquad S_{2n+1}=S_n+S_{n+1},}
e $CW(n)=S_n/S_{n+1}$, $n\ge1$, que enumera $\Q_{>0}$ sem repetições.
\begin{lema}
Para todo $q\ge1$, $CW(2^{q-1})=1/q$.
\end{lema}
\begin{proof}
Da recorrência par, $S_{2^m}=S_1=1$. Também $S_{2^m+1}=m+1$: o caso $m=0$ é $S_2=1$ e
\[
S_{2^{m+1}+1}=S_{2(2^m)+1}=S_{2^m}+S_{2^m+1}=1+(m+1).
\]
Tomando $m=q-1$, resulta $CW(2^{q-1})=1/q$.
\end{proof}

O racional $1/q$ tem altura $q$, mas índice $2^{q-1}$; a escrita binária do índice exige $q$ bits, contra os $O(\log q)$ do lema~2.4. A ordenação por altura controla o endereço em função de $\heightq(r)$. A seção~7 determina o comprimento mínimo no pior caso por contagem e constrói uma enumeração que o realiza.
\figura{01_calkin_wilf}{.94}{Primeiros níveis da árvore de Calkin--Wilf. O ramo esquerdo contém $1/q$ no índice $2^{q-1}$; a ordenação por altura fornece índices polinomiais na altura.}{cw}
